\documentclass[11pt,letterpaper]{moderncv}
\usepackage{multiaudience}
\usepackage{geometry}

\SetNewAudience{industry}
\SetNewAudience{academia}
\moderncvstyle{banking}
\moderncvcolor{blue}

\usepackage{xspace}
\usepackage{fontspec}
\usepackage{etoolbox}
\usepackage{enumitem}

\def\CurrentAudience{industry} % industry, academia, or policy

\SetNewAudience{industry}
\SetNewAudience{academia}
\SetNewAudience{policy}

\csdef{resumeconfig@industry}{%
  \def\ResumeMarginWidth{.70in}%
  \def\ResumeMarginHeight{.80in}%
  % TODO: Consider increasing bottom top margin, removing quote h
  % \def\ResumeHeadWidth{0.90\textwidth}% etc
}

\csdef{resumeconfig@academia}{%
  % Make resume margin bigger if I can afford it!!! 
  \def\ResumeMargin{0.50in}%
  % here since it has to be in the preamble
  \quote{Research Interests: programming languages, software systems, incremental computation}
}

% EXAMPLE -- unused
\csdef{resumeconfig@policy}{%
  \def\ResumeMargin{0.85in}%
}

\ifcsdef{resumeconfig@\CurrentAudience}
  {\csuse{resumeconfig@\CurrentAudience}}
  {%
    \PackageError{resume}
      {Unknown audience `\CurrentAudience'}
      {Valid audiences are industry, academia, and policy.}%
  }

\ifcsdef{ResumeMargin}
  {\geometry{margin=\ResumeMargin}}
  {%
    \ifcsdef{ResumeMarginWidth}
      {\geometry{
         left=\ResumeMarginWidth,
         right=\ResumeMarginWidth,
         top=\ResumeMarginHeight,
         bottom=\ResumeMarginHeight
       }}
      {%
        \PackageError{resume}
          {No resume margins defined}
          {Define either \string\ResumeMargin\space or both
           \string\ResumeMarginWidth\space and
           \string\ResumeMarginHeight.}%
      }%
  }
  
\newcommand{\entry}[4]{%
  \begin{tabular*}{\maincolumnwidth}[t]{%
    l@{\extracolsep{\fill}}r%
  }%
    \textbf{#1} & \textit{#2}%
    \ifstrempty{#3}{}{%
      \\
      \multicolumn{2}{@{}l@{}}{%
        \makebox[\maincolumnwidth][l]{#3}%
      }%
    }%
  \end{tabular*}%
  \ifstrempty{#4}{}{%
    \\%
    \begin{minipage}[b]{\maincolumnwidth}%
      #4%
    \end{minipage}%
  }%
  \par%
  \vspace{-\dp\strutbox}%
  \addvspace{.5em}%}
}

\makeatletter

\RenewDocumentCommand{\section}{sm}{%
  \par\addvspace{2ex}% changed from 2.5ex
  \phantomsection{}%
  \addcontentsline{toc}{section}{#2}%
  \if@left\else\if@fullrules\else\if@mixedrules\else%
    \sectionrule\fi\fi\fi%
  \strut\sectionstyle{#2}%
  \if@fullrules%
    \sectionrule%
  \else\if@mixedrules%
    \sectionrule%
  \else\if@right\else%
    \sectionrule\fi\fi\fi%
  \par\nobreak\addvspace{1ex}\@afterheading%
}

\makeatother


\newcommand{\pync}{{\scshape Pync}\xspace}

\newenvironment{presentationlist}{%
  \begin{itemize}[
    label={},
    leftmargin=0pt,
    labelsep=0pt,
    itemindent=0pt,
    topsep=0pt,
    itemsep=0.4ex,
    parsep=0pt,
    partopsep=0pt,
    before=\vspace{0.5ex}
  ]%
}{%
  \end{itemize}%
}

\newcommand{\presentationitem}[2]{%
  \item
  \noindent
  \begin{minipage}[t]{0.76\linewidth}%
    #1%
  \end{minipage}%
  \hfill
  \begin{minipage}[t]{0.22\linewidth}%
    \raggedleft\textit{#2}%
  \end{minipage}%
}


% for the quote
\patchcmd
  {\recomputecvheadlengths}
  {0.65\textwidth}
  {0.9\textwidth} % ideally would be .85
  {}
  {}
% Ideally would be 1em
\patchcmd{\makecvhead}{\\[2.5em]}{\\[.85em]}{}{}
\patchcmd{\makecvhead}{\\[2.5em]}{\\[.85em]}{}{} % has to be duplicated for some reason
% TODO: Fix.
\patchcmd{\section}
  {\addvspace{2.5ex}}
  {\addvspace{1.8ex}}
  {}
  {\PackageWarning{resume}{Could not patch section spacing}}

% Personal Information
\name{Bolun}{Thompson}
\email{bolunthompson@ucla.edu}
\phone{(858) 999-9567}
\homepage{bolun.dev}
\social[github]{bolunthompson}
\social[linkedin]{bolunthompson}


\begin{document}

\makecvtitle\section{Education}
\entry
  {B.S. in Computer Science}
  {Expected Graduation: 06/2028}
  {University of California, Los Angeles (UCLA) \hfill GPA: 3.87}
  {Coursework: Programming Languages, Operating Systems, Algorithms,
   Distributed Systems, Security}
   
\cvitem[0pt]{Program}{Oregon Programming Languages Summer School; studied Rocq and separation logic \quad \hfill \textit{06/2026}}
\vspace{-\dp\strutbox}% done in the entry, but not here since cvitem ends

\begin{shownto}{academia}
\section{Research Experience}
\end{shownto}

\begin{shownto}{industry}
\section{Experience}
\end{shownto}

\entry
  {Undergraduate Researcher}
  {02/2025--Present}
  {Programmable Software Systems Lab, UCLA, mentored by Prof. Konstantinos Kallas}
  {%
    \begin{itemize}[leftmargin=*,topsep=8pt,itemsep=4pt,parsep=0pt]
      \item Designed and implemented \pync, an $11$k line incremental execution system for Python that tracks function dependencies and effects to reuse prior computations across script executions. Achieved up to $10\times$ speedups when re-executing eight real-world data science workloads.
      \begin{shownto}{industry} Presented work to three computer science workshops.\space
      \end{shownto} Manuscript in preparation.

      \item Built and evaluated a 700 line Python front-end for hS, a speculative parallelization system for shell scripts. Uses AST transformations to identify and extract subprocess calls while preserving script semantics.
      Enabled hS to achieve up to $8\times$ speedups on five subprocess-heavy real-world Python workloads.

    \begin{shownto}{industry}  
      \cvitem{Publication}{\textit{hS: Speculative Script Reordering at Subprocess Granularity}. Co-author. \hfill OSDI '26}
    \end{shownto}
    \end{itemize}%
  }


\begin{shownto}{academia}
\section{Publications}

\entry
  {hS: Speculative Script Reordering at Subprocess Granularity.}
  {OSDI~'26}
  {}
  {Georgios Liargkovas, Di Jin, Tianyu (Ezri) Zhu, Dan Liu,
   \textbf{A. Bolun Thompson}, Anirudh Narsipur, Seong-Heon Jung,
   Siddhartha Prasad, Diomidis Spinellis, Michael Greenberg,
   Konstantinos Kallas, Nikos Vasilakis.}

% \entry
%  {\pync: Function-Level Incremental Execution for Python Scripts.}
%  {2026}
%  {}
%  {A. Bolun Thompson.
%   PLDI '26 Student Research Competition.}
\end{shownto}

\begin{shownto}{academia}
\section{Honors, Awards and Service}
\end{shownto}

\begin{shownto}{industry}
\section{Honors and Awards}
\end{shownto}

\entry
  {First Place, undergraduate division --- PLDI '26 Student Research Competition}
  {06/2026}
  {}
  {Presented \pync, an automatic incremental execution system
   for Python.}
\begin{shownto}{academia}
\entry
  {Student Volunteer --- PLDI '26}
  {06/2026}
  {}
  {}
\end{shownto}
\section{Open-Source Contributions and Projects}
% TODO: Post on GH, review if asked about before interview.
% TODO Post on github.
\entry
  {Secure Cloud Storage}
  {03/2026}
  {}
  {Built a 2k-line secure cloud storage client in Go on top of an untrusted server, using AES/HMAC primitives.
  Maintains confidentiality and integrity. Supports multiple users, file sharing, and revocation.
  }
\begin{shownto}{industry}
% Maybe mention security project, DS project, review these implementations?
\entry{Paper Plane}{09/2025--11/2025}{}{
Co-developed a full-stack flight-log web application that enables pilots to manage their flight history. Built the REST API backend with Express.js and TypeScript, including user authentication; designed the Prisma/SQLite data model; and wrote backend integration tests executed in CI.}

\end{shownto}
\entry{PaSh}{11/2024--02/2025}{}{
Contributed N pull requests to PaSh, a system for automatically parallelizing shell scripts. Added Ubuntu 24.04 compatibility and completed a new Bash front-end by implementing word-expansion.
% TODO: Add somethign about impact.
}

\begin{shownto}{academia}
\section{Research Presentations}

\entry
  {\pync: Python Incremental Execution}
  {}
  {}
  {%
  \begin{presentationlist}    
    \presentationitem
      {Talk, Undergraduate Research Showcase, UCLA}
      {05/2025, 05/2026}
    \presentationitem
      {Invited Talk, SoCal PLs Day, University of Southern California}
      {02/2026}
    \presentationitem
      {Talk, 5th PaSh Workshop, Brown University}
      {09/2025}
  \end{presentationlist}
  }
\end{shownto} 

\section{Skills}

\cvitem[0.2em]{Languages}
  {Python, C, C++, JavaScript, Rust, Go, Shell, OCaml, Rocq, Dafny, Prolog}

\cvitem{Tools}
  {Git, Docker, GitHub Actions, Linux, Pytest, \LaTeX}
\end{document}